\setcounter{chapter}{-1}
\chapter{Introduction}
\textbf{BLOCKS OF THE COURSE}:
\begin{itemize}
    \item Measure (measurable functions) and integral (integrable function)
    \item MTC, DCT
    \item Fubini Thm
    \item $L^{p}$ space


\end{itemize}
\begin{eg}
    SUPER IMPORTANT EXAMPLE!! Let $f_n(x) = n^2x^{n}(1-x) (0 \le x \le 1)$. Then $\lim_{n \to \infty} f_n(x) = 0$ for all $x \in  [0, 1]$, but $\lim_{n \to \infty} \int_{0}^{1} f_n(x)dx = 1$. This example arises in any reasonable theory.
\end{eg}
We want theorems of the form $(f_n)$ is a sequence of integrable functions, $f_n(x) \to  f(x) $ for each $x$, and [\textit{additional assumptions}]. Then $f$ is integrable and $\int f(x) dx = \lim_{n \to \infty} \int f_n(x) dx$. 
\\
\textbf{Monotone Convergence Theorem, Dominated Convergence Theorem.}
\\
Crucial ideas for Lebesgue's construction:
\begin{itemize}
    \item Instead of using integrals to define lengths of sets, define length of a set directly; then define integrals.
    \item Instead of partitioning the $x$-axis into intervals and using step functions, partition the $y$-axis into intervals and consider corresponding "simple" functions.
\end{itemize}
